% !TeX root = ../navier-stokes-manuscript.tex
\clearpage
\section{The Estimates and Global Smoothness}\label{sec:TheMissingEstimateAndGlobalSmoothness}

\subsection{Every finite estimate is a coordinate}\label{sub:FromFiniteRectanglesToClassicalEstimates}

Fix a temporal order \(r\in\Nzero\) and a spatial order \(m\in\Nzero\).
Choose one rectangle with \(N\geq r\) and middle height \(M=m\).
Its adjacent rows give
\[
  V_r\in L^2(I_*;H^{m+1}),
  \qquad
  \partial_tV_r=V_{r+1}\in L^2(I_*;H^{m-1}).
\]
The trace identity \eqref{eq:BodyAdjacentTraceIdentity} gives
\begin{equation}\label{eq:BodyUniformMixedSobolevCoordinate}
  \boxed{
  M_{r,m}
  :=\sup_{0\leq t\leq T_*}
  \norm{\partial_t^ru(t)}_{H^m}
  <\infty.}
\end{equation}
The pair \((r,m)\) was arbitrary.
Every fixed temporal response is uniformly controlled through \(T_*\) in every selected finite Sobolev space.

That one statement generates the ordinary finite-order estimates.
Let \(k\in\Nzero\), let \(2\leq p\leq\infty\), and choose an integer \(m\) with
\begin{equation}\label{eq:SobolevProjectionThreshold}
  m>k+\frac32-\frac3p.
\end{equation}
Spatial Sobolev embedding gives
\begin{equation}\label{eq:MasterSpatialProjection}
  \norm{\nabla^k\partial_t^ru(t)}_{L^p}
  \leq C_{k,m,p}\norm{\partial_t^ru(t)}_{H^m}.
\end{equation}
Using \eqref{eq:BodyUniformMixedSobolevCoordinate},
\begin{equation}\label{eq:MasterUniformMixedEstimate}
  \nabla^k\partial_t^ru
  \in L^\infty(I_*;L^p(\R^3)).
\end{equation}
Since \(T_*\) is finite under the standing contradiction hypothesis,
\begin{equation}\label{eq:MasterFiniteTimeMixedEstimate}
  \nabla^k\partial_t^ru
  \in L^q(I_*;L^p(\R^3))
  \qquad(1\leq q\leq\infty).
\end{equation}
Each requested norm uses one finite rectangle and its own finite constant.
No sum over all derivative orders and no common analytic radius is being asserted.

The reversal from Section~3 is now an estimate rather than a pattern in an identity.
For any fixed temporal row \(V_r\), the singularity argument gave
\[
  \sup_{0<t<T_*}\norm{V_r(t)}_3=\infty.
\]
The coordinate \((r,m)=(r,1)\) gives instead
\[
  \sup_{0\leq t\leq T_*}\norm{V_r(t)}_3
  \leq C M_{r,1}<\infty.
\]
The temporal and spatial directions have not been identified with one another.
They are controlled simultaneously inside the same history, and the same row cannot satisfy both conclusions.

The lowest rectangles now return the argument to the critical height and the vortex that opened the Tower.

\subsection{The critical-height roof}\label{sub:ClassicalContinuationCriteriaInsideFiniteRectangles}

The first rectangle that reaches the critical endpoint is \(\cR_{0,1}\).
It gives
\begin{equation}\label{eq:LowRectangleZeroOne}
  u\in L^2(I_*;H^2),
  \qquad
  u_t\in L^2(I_*;L^2),
  \qquad
  u\in C([0,T_*];H^1).
\end{equation}
Hence
\begin{equation}\label{eq:EndpointCriterionFromLowRectangle}
  \sup_{0\leq t\leq T_*}\norm{u(t)}_3
  \leq C\sup_{0\leq t\leq T_*}\norm{u(t)}_{H^1}
  <\infty.
\end{equation}
This is the same \(L^3\) norm that a finite maximal time forced to escape in \eqref{eq:FiniteTimeLThreeEscape}.
The lowest critical projection of the intersection already blocks that escape.

It also gives the critical height a finite roof:
\begin{equation}\label{eq:BodyCriticalHeightRoof}
  H_{\max}
  :=\sup_{0\leq t\leq T_*}H(t)
  \leq\frac12\sup_{0\leq t\leq T_*}\norm{u(t)}_{H^1}^2
  <\infty.
\end{equation}
Only finitely many of the first-passage levels \(L_k=2^kH_0\) fit below this roof.
If \(N_{\mathrm{levels}}\) denotes the number reached, then
\begin{equation}\label{eq:BodyFiniteZenoOctaveCount}
  N_{\mathrm{levels}}
  \leq
  \begin{cases}
    0,&H_{\max}<H_0,\\[0.25em]
    1+\left\lfloor\log_2(H_{\max}/H_0)\right\rfloor,
      &H_{\max}\geq H_0.
  \end{cases}
\end{equation}
No estimate on the duration of an individual crossing is needed.
The infinite Zeno ladder fails because the height cannot cross beyond a finite roof.

The next low rectangle, \(\cR_{0,2}\), records the critical-height dynamics as well as its ceiling:
\begin{equation}\label{eq:LowRectangleZeroTwo}
  u\in L^2(I_*;H^3),
  \qquad
  u_t\in L^2(I_*;H^1),
  \qquad
  u\in C([0,T_*];H^2).
\end{equation}
The first two memberships give
\[
  H'(t)
  =\operatorname{Re}
  \pair{\Lambda^{1/2}u(t)}{\Lambda^{1/2}u_t(t)}_{L^2}
  \in L^1(I_*).
\]
Also \(E_{3/2}\in L^1(I_*)\), so the first completed row
\[
  \mathcal P_{1/2}=H'+2\nu E_{3/2}
\]
shows that the signed critical production is integrable on the whole candidate interval.
The height can rise and fall sharply, but it cannot accumulate an infinite net rise before \(T_*\).

Sobolev embedding gives the finite stretching action
\begin{equation}\label{eq:FiniteLipschitzVorticityStrainAction}
  \int_0^{T_*}
  \left(
    \norm{\nabla u(t)}_\infty
    +\norm{\omega(t)}_\infty
    +\norm{S(t)}_\infty
  \right)dt
  <\infty.
\end{equation}
This is the estimate needed to return from the global critical-height obstruction to the two physical readings of the balanced vortex.

\subsection{The two vortex readings}\label{sub:ClosingCriticalHeightZenoAndTheVortex}

Return first to a core made from the same material labels.
If \(X(a,t)\) and \(X(b,t)\) are two trajectories in that core, then
\[
  \frac d{dt}|X(a,t)-X(b,t)|
  \geq
  -\norm{\nabla u(t)}_\infty|X(a,t)-X(b,t)|.
\]
The finite action in \eqref{eq:FiniteLipschitzVorticityStrainAction} gives
\begin{equation}\label{eq:BodyMaterialCoreLowerBound}
  |X(a,t)-X(b,t)|
  \geq
  |a-b|
  \exp\!\left(
    -\int_0^t\norm{\nabla u(s)}_\infty\,ds
  \right)
  >0
\end{equation}
for every \(t\leq T_*\).
A persistent material core containing two labels with positive initial separation cannot collapse those labels to one point in finite time.

Perfect centring does not open a loophole.
The finite coordinate \((r,m)=(0,3)\) in \eqref{eq:BodyUniformMixedSobolevCoordinate} gives
\begin{equation}\label{eq:BodyUniformLipschitzBound}
  L_*:=\sup_{0\leq t\leq T_*}\norm{\nabla u(t)}_\infty<\infty.
\end{equation}
At a chosen centre \(x_c(t)\),
\[
  |u(t,x)-u(t,x_c)|\leq L_*|x-x_c|.
\]
Let \(C_r(t)\) be a circle of radius \(r\) in a cross-section normal to the vortex axis.
The constant vector \(u(t,x_c)\) integrates to zero around the closed curve, so
\begin{equation}\label{eq:BodyCenteredCirculationBound}
  \left|
    \oint_{C_r(t)}u(t,x)\cdot d\ell
  \right|
  \leq2\pi L_*r^2
  \longrightarrow0
  \qquad(r\downarrow0).
\end{equation}
The relative tangential speed is \(O(r)\), and the circulation is \(O(r^2)\).
Exact balance at the centre does not preserve nonzero circulation at zero radius.

Now take the second reading, in which the visible core recruits new material as earlier labels leave it.
Particle separation no longer follows the changing membership of that core, so the test moves to the whole-field high-frequency tail.
For \(s>1/2\) and \(K\geq1\), define
\[
  H_{>K}(t)
  :=\frac12\int_{|\xi|>K}
  |\xi|\,|\widehat u(t,\xi)|^2\,d\xi.
\]
Then
\begin{equation}\label{eq:BodyUltravioletCriticalTail}
  H_{>K}(t)
  \leq
  \frac12K^{1-2s}\norm{u(t)}_{\dot H^s}^2.
\end{equation}
Choose an integer \(m\geq s\).
The coordinate \(M_{0,m}<\infty\) and the continuous embedding \(H^m\hookrightarrow\dot H^s\) give
\begin{equation}\label{eq:BodyUniformUltravioletVanishing}
  \lim_{K\to\infty}
  \sup_{0\leq t\leq T_*}H_{>K}(t)=0.
\end{equation}
Even when material is continually recruited, the complete velocity field cannot retain a fixed amount of critical content beyond every frequency.

These conclusions do not use wobble, fragmentation, or broken symmetry.
Persistent labels retain positive separation, exact centring carries vanishing circulation into a vanishing radius, and recruitment leaves a uniformly vanishing whole-field critical tail.
The perfectly balanced vortex has no remaining route to a singular zero-scale concentration at \(T_*\).

\subsection{No finite singular time}\label{sub:NoFiniteSingularTime}

\begin{theorem}[Whole-space existence and smoothness]\label{thm:MainGlobalRegularity}
For every \(\nu>0\) and every \(u_0\in\Ssigma\), the maximal classical solution of \eqref{eq:NavierStokesSystem} satisfies
\begin{equation}\label{eq:BodyGlobalSmoothnessConclusion}
  T_*=\infty,
  \qquad
  u,p\in C^\infty\!\bigl(\R^3\times[0,\infty)\bigr).
\end{equation}
For every \(t\geq0\),
\begin{equation}\label{eq:GlobalEnergyConclusion}
  \frac12\norm{u(t)}_2^2
  +\nu\int_0^t\norm{\nabla u(s)}_2^2\,ds
  =\frac12\norm{u_0}_2^2.
\end{equation}
\end{theorem}

\begin{proof}
Assume that \(T_*<\infty\).
The whole-terminal estimate gives every finite rectangle, and \eqref{eq:BodyProducedMixedIntersection} assembles their common coordinates.
Applying \eqref{eq:BodyUniformMixedSobolevCoordinate} with \(r=0\) gives one terminal velocity
\begin{equation}\label{eq:BodyCommonSmoothEndpoint}
  u_*:=u(T_*)
  \in\bigcap_{m=0}^{\infty}H^m(\R^3)
  =\Hinfty(\R^3).
\end{equation}
The traces at different Sobolev heights are limits of the same function and agree under the natural embeddings.

The normalized pressure recurrence and the differentiated momentum recurrence determine the complete left endpoint jet from \(u_*\).
Local well-posedness starting from \(u_*\) produces a smooth solution on \([T_*,T_*+\delta]\), and uniqueness joins it to the original solution across the endpoint \parencite{Kato1984}.
Appendix~\ref{app:CompatibleTerminalAssembly} gives the trace compatibility, pressure-jet passage, and smooth gluing in full.
The joined solution exists beyond \(T_*\), contradicting maximality.

Thus \(T_*=\infty\), which proves \eqref{eq:BodyGlobalSmoothnessConclusion} and the time requirement \eqref{eq:GlobalTimeTarget}.
The energy identity \eqref{eq:BodyEnergyIdentity} holds on every finite interval of the global solution and gives \eqref{eq:GlobalEnergyConclusion}, completing \eqref{eq:WholeSpaceRegularityTarget}.
\end{proof}

\clearpage
\section{Further Consequences and the Euler Boundary}\label{sec:FurtherConsequencesAndEulerBoundary}

\subsection{A finite sum across infinitely many scales}\label{sub:AFiniteSumAcrossScales}

The rectangle \(\cR_{0,2}\) controls one finite Sobolev row, yet that row sums the vorticity carried by every dyadic scale.
Let \((\Delta_j)_{j\geq-1}\) be an inhomogeneous Littlewood--Paley decomposition \parencite[Chapter~2]{BahouriCheminDanchin2011}.
For \(j\geq0\), the block \(\Delta_j\omega\) is concentrated near wavelength \(2^{-j}\), and \(\norm{\Delta_j\omega}_\infty\) measures the strongest rotation present in that band.
Define
\[
  \norm{\omega}_{B^0_{\infty,1}}
  :=\sum_{j\geq-1}\norm{\Delta_j\omega}_\infty.
\]
Bernstein's inequality gives, for \(j\geq0\),
\[
  \norm{\Delta_j\omega}_\infty
  \leq C2^{3j/2}\norm{\Delta_j\omega}_2.
\]
Insert \(2^{3j/2}=2^{-j/2}2^{2j}\), sum, and use Cauchy--Schwarz in the shell index:
\begin{align*}
  \sum_{j\geq0}\norm{\Delta_j\omega}_\infty
  &\leq
  C\left(\sum_{j\geq0}2^{-j}\right)^{1/2}
  \left(\sum_{j\geq0}2^{4j}\norm{\Delta_j\omega}_2^2\right)^{1/2}
  \\
  &\leq C\norm{\omega}_{H^2}
  \leq C\norm{u}_{H^3}.
\end{align*}
The low block \(\Delta_{-1}\omega\) is controlled by \(\norm{\omega}_2\), so
\begin{equation}\label{eq:VorticityBesovFromHThree}
  \norm{\omega}_{B^0_{\infty,1}}
  \leq C\norm{u}_{H^3}.
\end{equation}
Since \(\cR_{0,2}\) gives \(u\in L^2(I_*;H^3)\), Cauchy--Schwarz in time yields
\begin{equation}\label{eq:FiniteDyadicVorticityAction}
  \sum_{j\geq-1}
  \int_0^{T_*}\norm{\Delta_j\omega(t)}_\infty\,dt
  =\int_0^{T_*}\norm{\omega(t)}_{B^0_{\infty,1}}\,dt
  <\infty.
\end{equation}
This is a finite \(\ell^1\) value across infinitely many possible shells, not a claim that only finitely many shells can become active.
A recruited core may pass activity from one scale to the next, but the time-integrated shell suprema cannot form a nonsummable cascade on a finite interval.
The same low rectangle supplies \(L^2_tH^3_x\) for this shell action and \(C_tH^2_x\) for the uniform critical-tail estimate \eqref{eq:BodyUniformUltravioletVanishing}.
The conclusion is not an analytic or Gevrey estimate: each finite order has its own rectangle, with no weighted all-order sum asserted.

\subsection{The critical endpoint seen from the other side}\label{sub:ClassicalCriteriaInsideFiniteRectangles}

The \(L^3\) endpoint first entered as a necessary escape.
If \(T_*\) were finite, \eqref{eq:FiniteTimeLThreeEscape} would force \(\norm{u(t)}_3\) to become unbounded along the terminal approach.
The rectangle \(\cR_{0,1}\) returns to that same norm from the other direction.
Its adjacent pair gives \(u\in C([0,T_*];H^1)\), and hence
\[
  \sup_{0\leq t\leq T_*}\norm{u(t)}_3<\infty.
\]
Thus \(\cR_{0,1}\) supplies the hypothesis of the Escauriaza--Seregin--\v{S}ver\'ak endpoint continuation theorem \parencite{EscauriazaSereginSverak2003}.
The norm that defined the singular threat is already bounded by the first critical rectangle.

The upper edge of the same rectangle gives another historical reading:
\[
  u\in L^2(I_*;H^2)
  \hookrightarrow L^2(I_*;L^\infty).
\]
This is the endpoint pair \((p,q)=(2,\infty)\) on the Ladyzhenskaya--Prodi--Serrin line \(2/p+3/q=1\) \parencite{Prodi1959,Serrin1962,Serrin1963,Ladyzhenskaya1967}.
These are two familiar measurements of the same finite block.
Any continuation criterion asking for finitely many temporal and spatial derivatives can be read from a sufficiently large rectangle by \eqref{eq:MasterFiniteTimeMixedEstimate}, with all readings attached to one velocity and one normalized pressure history.

\subsection{Why the Euler Tower is different}\label{sub:WhyTheEulerTowerIsDifferent}

The mixed derivative jet can be written for Euler as well as for Navier--Stokes, but the spatial ascent that drives this paper is viscous.
To keep the two histories distinct, write
\[
  V_n^{\mathrm{NS}}:=\partial_t^nu^{\mathrm{NS}},
  \qquad
  Q_n^{\mathrm{NS}}:=\partial_t^np^{\mathrm{NS}},
\]
and define \(V_n^{\mathrm E},Q_n^{\mathrm E}\) from a separate Euler solution.
Their temporal recurrences are
\begin{align}
  V_{n+1}^{\mathrm{NS}}
  &+
  \sum_{a=0}^{n}\binom{n}{a}
  (V_a^{\mathrm{NS}}\cdot\nabla)V_{n-a}^{\mathrm{NS}}
  +\nabla Q_n^{\mathrm{NS}}
  =\nu\Delta V_n^{\mathrm{NS}},
  \label{eq:NSTemporalRecurrenceComparison}
  \\
  V_{n+1}^{\mathrm E}
  &+
  \sum_{a=0}^{n}\binom{n}{a}
  (V_a^{\mathrm E}\cdot\nabla)V_{n-a}^{\mathrm E}
  +\nabla Q_n^{\mathrm E}
  =0.
  \label{eq:EulerTemporalRecurrenceComparison}
\end{align}
Both equations retain the differentiated transport and the nonlocal pressure response.
Only Navier--Stokes has the linear path from temporal row \(n\) to two additional spatial derivatives of that same row.

For the two separate histories, set
\[
  E_s^{\mathrm{NS}}:=\frac12\norm{\Lambda^su^{\mathrm{NS}}}_2^2,
  \qquad
  E_s^{\mathrm E}:=\frac12\norm{\Lambda^su^{\mathrm E}}_2^2,
\]
with corresponding transfer terms \(\mathcal P_s^{\mathrm{NS}}\) and \(\mathcal P_s^{\mathrm E}\).
Then
\begin{equation}\label{eq:NSEulerEnergyComparison}
  (E_s^{\mathrm{NS}})'
  =\mathcal P_s^{\mathrm{NS}}-2\nu E_{s+1}^{\mathrm{NS}},
  \qquad
  (E_s^{\mathrm E})'
  =\mathcal P_s^{\mathrm E}.
\end{equation}
For Navier--Stokes, repeated differentiation and division by \((-2\nu)^n\) isolate \(E_{n+1/2}^{\mathrm{NS}}\) in \eqref{eq:BodyCompletedCriticalRow}.
Euler retains the nonlinear mixing of time and space, but it has no signed coercive rung that isolates the next Sobolev height in this way.

An all-orders compatible intersection for an Euler history would still be a powerful conditional container.
It would imply every fixed Sobolev bound and, in particular, the Beale--Kato--Majda condition \(\int_0^T\|\omega^{\mathrm E}(t)\|_\infty\,dt<\infty\) \parencite{BealeKatoMajda1984}.
The fixed-viscosity theorem used here does not supply that Euler intersection, and its constants are not claimed to remain bounded as \(\nu\downarrow0\).
The comparison marks the exact boundary: the mixed jet survives without viscosity; the datum-generated viscous spatial ladder does not.
